\section{Local estimates without synchronized updates}\label{sec:local}

We work in one sector. Every correction ring toggles the boundary of
one plaquette, conditionally on its two incoming edges, and a sector
noise mark toggles the same boundary without a test. Write a vertex
as occupied when it is incident to at least one marked edge. The
following deterministic facts refer to the actual sequence of rings.

\subsection{Confinement and a clearance schedule}

\begin{lemma}[Geometric confinement]\label{lem:confinement}
For a finite marked edge set on \(\Z^4\) evolving with no noise under
any sequence of upper-corner correction attempts, the coordinate box
of occupied vertices cannot enlarge and their maximum height cannot
increase. Each vertex degree modulo two is preserved. The confinement
assertions hold even if the starting degrees are odd.
\end{lemma}
\begin{proof}
An enabled face at \(v\) has its tested edges present. Consequently
\(v,v-e_i,v-e_j\) are already occupied. The remaining corner
\(v-e_i-e_j\) lies in their coordinate box and has smaller height than
\(v\). All vertices on potentially new edges are therefore in the
existing box and below or at its maximum height. Toggling a plaquette
boundary changes the degree of each of its four vertices by an even
integer modulo two. Induction over the attempts proves every
assertion. Ineffective attempts change nothing.
\end{proof}

Fix an integer \(H\) and a vertex \(x\) with
\(\operatorname{ht}(x)\le H\). The northeast scheduling cone is
\[
 \mathcal C(x,H)=\{x+a:a\in\Z_{\ge0}^4,\
                         \operatorname{ht}(x+a)\le H\}.
\]
At each vertex \(v\), let \(F_v(s)\) be the time of completion of the
following subsequence of its actual clocks after time \(s\): wait for
the six pairs \(i<j\) in lexicographic order, and repeat that list.
Other rings, including rings of these clocks while waiting for a
different pair, are allowed. For a fixed realization each next-ring
function, and hence \(F_v\), is nondecreasing.

Set \(S(v)=0\) above height \(H\), and recursively set
\begin{equation}\label{eq:schedule}
 S(v)=F_v\left(\max_{1\le i\le4}S(v+e_i)\right)
       \qquad(v\in\mathcal C(x,H)).
\end{equation}
Times here are measured relative to the beginning of a clean interval.

\begin{lemma}[Poisson clearance]\label{lem:clearance}
Suppose a clean process has initial maximum occupied height at most
\(H\), even degree at every vertex of \(\mathcal C(x,H)\), and all
faces incident to that cone available for correction. This includes
a finite cut if the cone and its incident faces lie in its interior.
After time \(S(x)\), the vertex \(x\) has no incident marked edges,
permanently. With independent rate-one clocks and
\(m=H-\operatorname{ht}(x)+1\),
\begin{equation}\label{eq:clearance-tail}
 \Prb\{S(x)>t\}\le
 4^m2^{12m}e^{-t/2}\qquad(t\ge0).
\end{equation}
The event \(S(x)\le t\) reads only the clocks up to \(t\).
\end{lemma}
\begin{proof}
An incoming edge at a vertex above height \(H\) is always absent by
Lemma~\ref{lem:confinement}. Suppose the four vertices \(v+e_i\)
have permanently lost their incoming edges by time \(s\). The four
outgoing edges from \(v\) are then absent. An incoming edge at \(v\)
can be affected only by a face with upper corner \(v\), or one of
those four upper neighbours. A face based at an upper neighbour
cannot subsequently be enabled, because its incoming edges are
absent.

The number of incoming marked edges at \(v\) is now even by parity.
Every enabled face based at \(v\) removes two such edges; none can
create an incoming edge there. In one full pass through the six
pairs, any fixed remaining pair either rings while still present,
or some removal has already happened. Thus a pass removes at least
two whenever at least two remain. There are at most four incoming
edges. Two passes remove them all. Thereafter no tested pair can
be enabled at \(v\), and the upper neighbours cannot recreate an
incoming edge. Descending induction on height proves the deterministic
claim using \eqref{eq:schedule}. Boundary parity in a cut is unnecessary
away from the indicated cone.

For the probability bound, monotonicity of \(F_v\) implies
\(F_v(\max_i s_i)=\max_iF_v(s_i)\). Expanding
\eqref{eq:schedule} therefore expresses \(S(x)\) as the maximum of
serial completions along northeast paths of \(m\) vertices, read from
the largest height downwards. There are at most \(4^m\) such paths.
Along any specified path the upper vertices are distinct. Its serial
completion time is a sum of \(12m\) independent rate-one exponentials.
Indeed the start time for the clocks at each successive vertex depends
only on clocks at distinct earlier vertices, and the next waiting
times at that vertex have the exponential law by independent increments
and the strong Markov property. For a rate-one exponential \(W\),
\(\E e^{W/2}=2\). The exponential Markov inequality followed by a path
union bound gives \eqref{eq:clearance-tail}.

Finally the recursive test for completion by \(t\) either finds all
required subsequence rings by \(t\), or fails to do so. It never needs
to inspect a ring after \(t\). This proves its measurability.
\end{proof}

\begin{corollary}[An isolated finite syndrome]\label{cor:island}
An even syndrome initially confined to an integer box of maximum
coordinate side length \(d\), with no noise, becomes empty by time
\(t\) except with probability at most
\[
 \min\{1,(d+1)^4(4\cdot2^{12})^{4d+1}e^{-t/2}\}.
\]
Every correction stays in that box. If the initial error chain is
supported in the box, its final closed residual is a boundary.
The same assertions hold for a torus box whose one-cell collar lifts
injectively, with independent torus clocks restricted to that box.
\end{corollary}
\begin{proof}
Choose \(H\) to be the initial maximum height. For any vertex in the
box, \(H-\operatorname{ht}(x)\le4d\). Apply
Lemma~\ref{lem:clearance} and sum over its at most \((d+1)^4\)
vertices. Confinement gives the support assertion. A closed 2-chain
in a cubical box bounds there; an explicit chain contraction is
provided in Section~\ref{sec:decoder}.

For the torus assertion transfer the clocks in the box and its collar
to its integer lift, and add independent clocks outside that collar.
Confinement makes every enabled face stay inside the original box,
so projection of this infinite-lattice construction agrees with the
confined torus trajectory. The auxiliary scheduling cone may extend
beyond one period. Its additional clocks are independent extension
clocks, rather than repeated torus clocks. The estimate thus uses
distinct independent vertices along each path and remains valid.
\end{proof}

The last distinction is useful for isolated components but will not
justify arbitrarily deep lifts of interacting torus histories. The
short-window construction below is what guarantees distinct physical
defects in a contour.

\subsection{Spatial blocks and the local test}

Fix an integer \(B\ge10\). For \(L\) sufficiently large define
\begin{equation}\label{eq:block-parameters}
 M=\lfloor L/B\rfloor,\quad b=L/M,\quad
 T=1000B,\quad m_f=1000T,\quad R=2000100.
\end{equation}
The coarse lattice is \(\Z_M^4\). In a consistent lift its blocks are
\[
 Q_z=\{v\in\Z^4:bz_i\le v_i<b(z_i+1),\quad1\le i\le4\}.
\]
Here \(B\le b<2B\). Translation by \(M\) in a block coordinate translates
the physical coordinate by exactly \(L\). Throughout this section
assume \(M>10(R+5)\). A bin \((z,k)\) refers to the physical slab
\(I_k=((k-1)T,kT]\). Lift the blocks with
\(\|d\|_\infty\le R\) relative to \(z\), and let
\[
 U_z=\bigcup_{\|d\|_\infty\le R}Q_{z+d}.
\]
To compare with the device, retain the initial syndrome edges whose
two endpoints lie in \(U_z\), delete the other edges, and run the
clean dynamics using only faces entirely in \(U_z\). Its clocks are
the physical clocks of these faces.

\begin{definition}[Bad bin]\label{def:bad}
The indicator \(\Bad(z,k)\) equals zero when all three conditions below
hold, and equals one otherwise:
\begin{enumerate}
\item No sector noise mark in \(I_k\) occurs on a face with any vertex
in \(U_z\).
\item There is no strictly chronological sequence of \(m_f\) correction
rings on successively vertex-intersecting faces, ending on a face
within physical sup distance two of \(Q_z\). Count such sequences
in the lifted neighbourhood of \(Q_z\).
\item For
\(H_z=\lfloor b(z_1+z_2+z_3+z_4+4)\rfloor\), every \(x\in Q_z\)
has \(S(x)\le T\), with the schedule of
\eqref{eq:schedule} started at the beginning of \(I_k\).
\end{enumerate}
\end{definition}

Each sequence in condition 2 lies within distance \(m_f+4\) of
\(Q_z\), and hence in \(U_z\). Each cone in condition 3 has depth
at most \(8B+2\), and its incident faces are interior to \(U_z\).
These tests depend solely on marks, not on the starting syndrome.

\begin{lemma}[Probability and range of dependence]\label{lem:bad-prob}
For every \(q>0\) there exist an integer \(B\ge10\) and
\(\eta_0>0\), independent of \(L\), such that for every allowed set of
rates bounded by \(\eta_0\), all \(z,k\), and all past histories up to
\((k-1)T\),
\[
 \Prb(\Bad(z,k)=1\mid\text{past})\le q.
\]
The test reads only face processes anchored in blocks at coarse
distance at most \(R+3\) from \(z\), during \(I_k\).
Disjoint sets of such basic increments are independent.
\end{lemma}
\begin{proof}
There are at most \(C_RB^4\) faces incident to \(U_z\). The sector
noise rate per face is at most \(2\eta\), so condition 1 fails with
probability at most \(2\eta C_RB^4T\).

A vertex of the four-dimensional cubical lattice is incident to
\(4\binom42=24\) faces. A face has four vertices, and so meets at
most 96 faces at vertices, including itself. There are at most
\(CB^4\) choices of a terminal face in condition 2. For each specified
list of \(m_f\) face labels, the expected number of ordered ring
tuples in a slab is \(T^{m_f}/m_f!\). This follows by integrating
the rate-one product intensity over the time simplex. If a face
label repeats, the factorial moment measure of its Poisson process
gives the same expression; ordered tuples use distinct times.
Consequently condition 2 fails with probability at most
\[
 CB^4\,\frac{(96T)^{m_f}}{m_f!}
 \le CB^4(96eT/m_f)^{m_f}.
\]
There are at most \((2B+1)^4\) target vertices in condition 3.
Lemma~\ref{lem:clearance} bounds its failure probability by
\[
 (2B+1)^4(4\cdot2^{12})^{8B+2}e^{-T/2}.
\]
Combining these estimates gives
\begin{equation}\label{eq:bad-cost}
 q_B(\eta):=2\eta C_RB^4T+
 CB^4(96e/1000)^{1000T}
 +(2B+1)^4(4\cdot2^{12})^{8B+2}e^{-T/2}.
\end{equation}
The last two terms tend to zero as \(B\) tends to infinity:
\(96e/1000<1\), and
\(8\log(4\cdot2^{12})-500<0\).
Choose \(B\) to make their sum at most \(q/2\), and then choose a
strictly positive \(\eta_0\) to make the noise term at most \(q/2\).

The face anchors needed for the three tests are in the stated
\(R+3\) block neighbourhood. All tests use increments within \(I_k\).
Those increments are independent of the entire past, even if the
initial syndrome was selected from it. The conditional assertion and
the independence statement follow.
\end{proof}

\subsection{Five predecessor regions}

Let
\begin{equation}\label{eq:regions}
 \begin{split}
 \Delta&=\{d\in\Z^4:\|d\|_\infty\le R\},\\
 A_i&=\{d\in\Delta:d_i\le0\}\quad(1\le i\le4),\\
 A_5&=\{d\in\Delta:d_1+d_2+d_3+d_4\ge1\}.
 \end{split}
\end{equation}
All five sets are nonempty and \(\bigcup_s A_s=\Delta\).

\begin{lemma}[Uniform predecessor implications]\label{lem:predecessor}
For every even starting syndrome and every realization with
\(\Bad(z,k)=0\), occupation of \(Q_z\) at time \(kT\) implies, for
each \(s=1,\ldots,5\), occupation at time \((k-1)T\) of some block
\(Q_{z+d}\) with \(d\in A_s\). Occupation of \(Q_z\) at any instant of
\(I_k\) implies occupation of some block of \(U_z\) at its start.
\end{lemma}
\begin{proof}
We justify agreement of the clean cut and the physical trajectory
near \(Q_z\) for the entire slab. Trace backwards any differing
incident edge. An interior correction ring can create a difference
only if a tested edge or the old output edge already differed.
Thus a discrepancy arriving from the cut boundary gives a
chronological chain of vertex-intersecting faces from that boundary
to a face within two of \(Q_z\). Initial interior edges agree, and
condition 1 excludes all relevant interior noise. In terms of upper
face anchors, consecutive intersecting faces differ by at most one
in each coordinate. Endpoint allowances are bounded by a fixed
number of cells. The cut boundary is at distance at least \(RB-2\),
which exceeds \(m_f+5\). Its final \(m_f\) interior rings would
contradict condition 2. This argument applies to a discrepancy at
any time of the slab. There are almost surely finitely many rings
on the finite torus in bounded time.

If no initial block with offset in \(A_i\) is occupied, every retained
marked endpoint has \(i\)th coordinate at least \(b(z_i+1)\).
Confinement prevents an endpoint from entering \(Q_z\), for the whole
slab. If no initial block in \(A_5\) is occupied, each retained
endpoint \(v\) satisfies
\[
 \operatorname{ht}(v)<b(z_1+z_2+z_3+z_4+4).
\]
Thus its maximum height is at most \(H_z\). All vertices of the
scheduling cones retain even initial degree, because they are
interior and no incident edge was deleted there. Condition 3 and
Lemma~\ref{lem:clearance} clear the target vertices by the slab end.
Finally, if \(U_z\) was initially empty, the clean cut stays empty
throughout the slab. Agreement with the physical trajectory proves
the two implications by contraposition.
\end{proof}

The monotonicity facts used here concern coordinate extrema, height
and the clock-completion functions. They require no ordering of
syndromes under inclusion and no synchronized correction rounds.
The bins are observations of the original marks; no part of the
generator is scheduled or reset at their boundaries.
